\sectionkicker{Source-backed technical notes}
\subsection*{生成から実行へ――Copilot、ReActとコード・行動・検索の接続: Technical Notes}
本欄は記事本文で圧縮した一次資料上の情報を、比較・再検証しやすい形で整理したものである。時系列、確認済みの主張、留意点、一次資料URLを掲載する。
\medskip
\subsection*{Theme at a glance}
\addcontentsline{toc}{subsection}{Theme at a glance}
\begin{center}
\footnotesize
\begin{tabularx}{\linewidth}{@{}p{0.20\linewidth}p{0.10\linewidth}p{0.14\linewidth}X@{}}
\toprule
資料 & 位置づけ & 種別 & 時系列 \\
\midrule
GitHub Copilot general availability & 主要資料 & PRODUCT & — \\
ReAct & 主要資料 & 論文 & — \\
AlphaCode & 補足資料 & モデル & — \\
Codex applications & 補足資料 & PRODUCT & — \\
Video PreTraining (VPT) & 補足資料 & Agent & — \\
Atlas & 補足資料 & 論文 & — \\
\bottomrule
\end{tabularx}
\end{center}
\normalsize

\begin{technicalnote}{GitHub Copilot general availability}{主要資料}
\begingroup
% reader-facing Technical Notes break policy
\widowpenalty=10000
\clubpenalty=10000
\displaywidowpenalty=10000
\begin{tabularx}{\linewidth}{@{}>{\bfseries}p{0.22\linewidth}X@{}}
組織 & GitHub \\
種別 & 製品 \\
時系列 & — \\
\end{tabularx}
\smallskip
{\bfseries 一次資料から整理したtechnical points}
\begin{itemize}[leftmargin=1.5em,itemsep=0.35em]
\item \textbf{一次情報で確認できる事実}: 一次資料で「GitHub Copilot general availability」が2022年の生成AI技術記録に含まれることを確認できる。技術内容、性能、アクセス、安全性に関する記述は、原著者・プロジェクトの主張として扱い、独立再現済みの結果とはみなさない。
\end{itemize}
\begin{minipage}{\linewidth}
% reader-facing Technical Notes generic boundary/source tail group
{\bfseries 読む際の境界}
\begin{itemize}[leftmargin=1.5em,itemsep=0.35em]
\item \textbf{分析上の留意点}: 一次資料で確認できる範囲の事実を記録しており、提供元・プロジェクト・著者による評価を独立再現済みの結果として扱わない。
\end{itemize}
\begin{samepage}
{\bfseries 一次資料}
\begingroup\sloppy
\Urlmuskip=0mu plus 2mu\relax
\begin{itemize}[leftmargin=1.5em,itemsep=0.25em]
\item {\scriptsize\url{https://github.blog/news-insights/product-news/github-copilot-is-generally-available-to-all-developers/}}
\end{itemize}
\endgroup
\end{samepage}
\end{minipage}
\endgroup
\end{technicalnote}
\medskip

\begin{technicalnote}{ReAct}{主要資料}
\begingroup
% reader-facing Technical Notes break policy
\widowpenalty=10000
\clubpenalty=10000
\displaywidowpenalty=10000
\begin{tabularx}{\linewidth}{@{}>{\bfseries}p{0.22\linewidth}X@{}}
組織 & authors \\
種別 & 論文 \\
時系列 & — \\
\end{tabularx}
\smallskip
{\bfseries 一次資料から整理したtechnical points}
\begin{itemize}[leftmargin=1.5em,itemsep=0.35em]
\item \textbf{Author claim}: 一次資料で「ReAct」が2022年の生成AI技術記録に含まれることを確認できる。技術内容、性能、アクセス、安全性に関する記述は、原著者・プロジェクトの主張として扱い、独立再現済みの結果とはみなさない。
\end{itemize}
\begin{minipage}{\linewidth}
% reader-facing Technical Notes generic boundary/source tail group
{\bfseries 読む際の境界}
\begin{itemize}[leftmargin=1.5em,itemsep=0.35em]
\item \textbf{分析上の留意点}: 一次資料で確認できる範囲の事実を記録しており、提供元・プロジェクト・著者による評価を独立再現済みの結果として扱わない。
\end{itemize}
\begin{samepage}
{\bfseries 一次資料}
\begingroup\sloppy
\Urlmuskip=0mu plus 2mu\relax
\begin{itemize}[leftmargin=1.5em,itemsep=0.25em]
\item {\scriptsize\url{http://arxiv.org/abs/2210.03629v3}}
\end{itemize}
\endgroup
\end{samepage}
\end{minipage}
\endgroup
\end{technicalnote}
\medskip

\begin{technicalnote}{AlphaCode}{補足資料}
\begingroup
% reader-facing Technical Notes break policy
\widowpenalty=10000
\clubpenalty=10000
\displaywidowpenalty=10000
\begin{tabularx}{\linewidth}{@{}>{\bfseries}p{0.22\linewidth}X@{}}
組織 & Google DeepMind \\
種別 & モデル \\
時系列 & — \\
\end{tabularx}
\smallskip
{\bfseries 一次資料から整理したtechnical points}
\begin{itemize}[leftmargin=1.5em,itemsep=0.35em]
\item \textbf{Author claim}: 一次資料で「AlphaCode」が2022年の生成AI技術記録に含まれることを確認できる。技術内容、性能、アクセス、安全性に関する記述は、原著者・プロジェクトの主張として扱い、独立再現済みの結果とはみなさない。
\end{itemize}
\begin{minipage}{\linewidth}
% reader-facing Technical Notes generic boundary/source tail group
{\bfseries 読む際の境界}
\begin{itemize}[leftmargin=1.5em,itemsep=0.35em]
\item \textbf{分析上の留意点}: 一次資料で確認できる範囲の事実を記録しており、提供元・プロジェクト・著者による評価を独立再現済みの結果として扱わない。
\end{itemize}
\begin{samepage}
{\bfseries 一次資料}
\begingroup\sloppy
\Urlmuskip=0mu plus 2mu\relax
\begin{itemize}[leftmargin=1.5em,itemsep=0.25em]
\item {\scriptsize\url{http://arxiv.org/abs/2203.07814v1}}
\item {\scriptsize\url{https://deepmind.google/blog/competitive-programming-with-alphacode/}}
\end{itemize}
\endgroup
\end{samepage}
\end{minipage}
\endgroup
\end{technicalnote}
\medskip

\begin{technicalnote}{Codex applications}{補足資料}
\begingroup
% reader-facing Technical Notes break policy
\widowpenalty=10000
\clubpenalty=10000
\displaywidowpenalty=10000
\begin{tabularx}{\linewidth}{@{}>{\bfseries}p{0.22\linewidth}X@{}}
組織 & OpenAI \\
種別 & 製品 \\
時系列 & — \\
\end{tabularx}
\smallskip
{\bfseries 一次資料から整理したtechnical points}
\begin{itemize}[leftmargin=1.5em,itemsep=0.35em]
\item \textbf{一次情報で確認できる事実}: 一次資料で「Codex applications」が2022年の生成AI技術記録に含まれることを確認できる。技術内容、性能、アクセス、安全性に関する記述は、原著者・プロジェクトの主張として扱い、独立再現済みの結果とはみなさない。
\end{itemize}
\begin{minipage}{\linewidth}
% reader-facing Technical Notes generic boundary/source tail group
{\bfseries 読む際の境界}
\begin{itemize}[leftmargin=1.5em,itemsep=0.35em]
\item \textbf{分析上の留意点}: 一次資料で確認できる範囲の事実を記録しており、提供元・プロジェクト・著者による評価を独立再現済みの結果として扱わない。
\end{itemize}
\begin{samepage}
{\bfseries 一次資料}
\begingroup\sloppy
\Urlmuskip=0mu plus 2mu\relax
\begin{itemize}[leftmargin=1.5em,itemsep=0.25em]
\item {\scriptsize\url{https://openai.com/index/codex-apps}}
\end{itemize}
\endgroup
\end{samepage}
\end{minipage}
\endgroup
\end{technicalnote}
\medskip

\begin{technicalnote}{Video PreTraining (VPT)}{補足資料}
\begingroup
% reader-facing Technical Notes break policy
\widowpenalty=10000
\clubpenalty=10000
\displaywidowpenalty=10000
\begin{tabularx}{\linewidth}{@{}>{\bfseries}p{0.22\linewidth}X@{}}
組織 & OpenAI \\
種別 & Agent \\
時系列 & — \\
\end{tabularx}
\smallskip
{\bfseries 一次資料から整理したtechnical points}
\begin{itemize}[leftmargin=1.5em,itemsep=0.35em]
\item \textbf{Author claim}: 一次資料で「Video PreTraining (VPT)」が2022年の生成AI技術記録に含まれることを確認できる。技術内容、性能、アクセス、安全性に関する記述は、原著者・プロジェクトの主張として扱い、独立再現済みの結果とはみなさない。
\end{itemize}
\begin{minipage}{\linewidth}
% reader-facing Technical Notes generic boundary/source tail group
{\bfseries 読む際の境界}
\begin{itemize}[leftmargin=1.5em,itemsep=0.35em]
\item \textbf{分析上の留意点}: 一次資料で確認できる範囲の事実を記録しており、提供元・プロジェクト・著者による評価を独立再現済みの結果として扱わない。
\end{itemize}
\begin{samepage}
{\bfseries 一次資料}
\begingroup\sloppy
\Urlmuskip=0mu plus 2mu\relax
\begin{itemize}[leftmargin=1.5em,itemsep=0.25em]
\item {\scriptsize\url{http://arxiv.org/abs/2206.11795v1}}
\item {\scriptsize\url{https://openai.com/index/vpt}}
\end{itemize}
\endgroup
\end{samepage}
\end{minipage}
\endgroup
\end{technicalnote}
\medskip

\begin{technicalnote}{Atlas}{補足資料}
\begingroup
% reader-facing Technical Notes break policy
\widowpenalty=10000
\clubpenalty=10000
\displaywidowpenalty=10000
\begin{tabularx}{\linewidth}{@{}>{\bfseries}p{0.22\linewidth}X@{}}
組織 & Meta AI / authors \\
種別 & 論文 \\
時系列 & — \\
\end{tabularx}
\smallskip
{\bfseries 一次資料から整理したtechnical points}
\begin{itemize}[leftmargin=1.5em,itemsep=0.35em]
\item \textbf{Author claim}: 一次資料で「Atlas」が2022年の生成AI技術記録に含まれることを確認できる。技術内容、性能、アクセス、安全性に関する記述は、原著者・プロジェクトの主張として扱い、独立再現済みの結果とはみなさない。
\end{itemize}
\begin{minipage}{\linewidth}
% reader-facing Technical Notes generic boundary/source tail group
{\bfseries 読む際の境界}
\begin{itemize}[leftmargin=1.5em,itemsep=0.35em]
\item \textbf{分析上の留意点}: 一次資料で確認できる範囲の事実を記録しており、提供元・プロジェクト・著者による評価を独立再現済みの結果として扱わない。
\end{itemize}
\begin{samepage}
{\bfseries 一次資料}
\begingroup\sloppy
\Urlmuskip=0mu plus 2mu\relax
\begin{itemize}[leftmargin=1.5em,itemsep=0.25em]
\item {\scriptsize\url{http://arxiv.org/abs/2208.03299v3}}
\end{itemize}
\endgroup
\end{samepage}
\end{minipage}
\endgroup
\end{technicalnote}
\medskip
